\chapter{Analyse et conception}

\section{Capture des besoins}
\subsection{Besoins fonctionnels}
\begin{longtable}{p{1.5cm}p{10cm}p{2cm}}
\caption{Besoins fonctionnels identifiés}\\
\toprule
\textbf{ID} & \textbf{Besoin} & \textbf{Module} \\
\midrule
BF1 & Créer un compte utilisateur avec nom, prénom, email, téléphone et mot de passe. & Auth \\
BF2 & Se connecter et conserver la session avec un token. & Auth \\
BF3 & Consulter et mettre à jour le profil. & Auth \\
BF4 & Consulter un tableau de bord utilisateur. & Frontend \\
BF5 & Consulter un tableau de bord administrateur avec nombre d'utilisateurs. & Admin \\
BF6 & Obtenir la météo actuelle et les prévisions. & Weather \\
BF7 & Recommander une culture selon localisation et données agro-climatiques. & Crops ML \\
BF8 & Personnaliser N, P, K, pH, température et humidité. & Crops ML \\
BF9 & Téléverser une image de plante pour diagnostic. & Analysis DL \\
BF10 & Obtenir top-3 prédictions, confiance et recommandations. & Analysis DL \\
BF11 & Echanger des messages internes. & Messaging \\
BF12 & Discuter avec un chatbot agricole. & Chatbot \\
BF13 & Consulter les actualités. & News \\
\bottomrule
\end{longtable}

\subsection{Besoins non fonctionnels}
\begin{itemize}[noitemsep]
  \item Sécurité : authentification par token DRF, permissions par endpoint, validation des entrées.
  \item Performance : chargement paresseux du CNN, cache LRU sur le géocodage des villes.
  \item Maintenabilité : applications Django séparées et services frontend isolés.
  \item Portabilité : rapport compatible Overleaf, application exécutable localement.
  \item Traçabilité : sauvegarde des analyses et conversations.
  \item Ergonomie : interfaces spécialisées par usage métier.
\end{itemize}

\subsection{Contraintes techniques}
Les contraintes détectées sont :
\begin{itemize}[noitemsep]
  \item backend Python/Django 4.2.7 ;
  \item frontend React 19 avec TypeScript ;
  \item base SQLite en environnement de développement ;
  \item dépendance à des APIs externes : Open-Meteo, NewsAPI, OpenAI ;
  \item modèle CNN volumineux \texttt{plant\_disease\_cnn.h5} ;
  \item modèles Random Forest sauvegardés en \texttt{.pkl}.
\end{itemize}

\section{Identification des acteurs}
\begin{table}[H]
\centering
\caption{Acteurs du système}
\begin{tabular}{p{3cm}p{10cm}}
\toprule
\textbf{Acteur} & \textbf{Description} \\
\midrule
Visiteur & Accède à l'accueil, à la connexion et à l'inscription. \\
Utilisateur & Accède aux services authentifiés : météo, recommandation, chatbot, analyse, messagerie. \\
Administrateur & Consulte le tableau de bord administrateur et la liste des utilisateurs via endpoint protégé. \\
API externe & Fournit météo, actualités ou réponses conversationnelles. \\
Modèle ML/DL & Produit des prédictions de culture ou de maladie. \\
\bottomrule
\end{tabular}
\end{table}

\section{Cas d'utilisation}
\subsection{Diagramme global}
\begin{figure}[H]
\centering
\caption{Diagramme global des cas d'utilisation}
\begin{lstlisting}[language=mermaid]
flowchart LR
    Visiteur((Visiteur))
    Utilisateur((Utilisateur))
    Admin((Administrateur))

    UC1[S'inscrire]
    UC2[Se connecter]
    UC3[Consulter meteo]
    UC4[Recommander culture]
    UC5[Analyser image plante]
    UC6[Utiliser chatbot]
    UC7[Gerer messages]
    UC8[Consulter dashboard admin]

    Visiteur --> UC1
    Visiteur --> UC2
    Utilisateur --> UC3
    Utilisateur --> UC4
    Utilisateur --> UC5
    Utilisateur --> UC6
    Utilisateur --> UC7
    Admin --> UC8
    Admin --> UC7
\end{lstlisting}
\end{figure}

\subsection{Descriptions textuelles}
\begin{longtable}{p{3cm}p{11cm}}
\caption{Descriptions des cas d'utilisation}\\
\toprule
\textbf{Cas} & \textbf{Description} \\
\midrule
Inscription & L'utilisateur renseigne ses informations. Le backend vérifie l'unicité de l'email et du nom d'utilisateur, compare les mots de passe et crée un compte. \\
Connexion & L'utilisateur fournit son nom d'utilisateur et son mot de passe. L'API authentifie et retourne un token stocké côté client. \\
Météo & L'utilisateur choisit une ville tunisienne. Le frontend appelle \texttt{/api/weather/current/} et \texttt{/api/weather/forecast/}. \\
Recommandation & L'utilisateur saisit une localisation et optionnellement des valeurs manuelles. Le backend récupère météo et sol, puis interroge Random Forest. \\
Analyse image & L'utilisateur envoie une image. L'API valide le fichier, prétraite l'image, appelle le CNN et sauvegarde le diagnostic. \\
Chatbot & L'utilisateur pose une question. Le backend envoie le contexte agricole et la question à l'API OpenAI puis sauvegarde la conversation. \\
Messagerie & L'utilisateur consulte ses messages reçus/envoyés, envoie un message, lit le détail ou supprime un message. \\
\bottomrule
\end{longtable}

\section{Modélisation UML}
\subsection{Diagramme de classes}
\begin{figure}[H]
\centering
\caption{Diagramme de classes principal}
\begin{lstlisting}[language=mermaid]
classDiagram
    class User {
      int id
      string username
      string email
      string role
      string phone
      image avatar
      datetime created_at
    }
    class Message {
      int id
      string subject
      text content
      boolean is_read
      datetime created_at
      datetime updated_at
    }
    class AnalysisLog {
      int id
      image image
      json result
      float confidence
      float processing_time
      datetime created_at
    }
    class ChatMessage {
      int id
      text message
      text response
      datetime created_at
    }
    User "1" --> "*" Message : sender
    User "1" --> "*" Message : recipient
    User "1" --> "*" AnalysisLog
    User "1" --> "*" ChatMessage
\end{lstlisting}
\end{figure}

\subsection{Diagramme de séquence : authentification}
\begin{figure}[H]
\centering
\caption{Séquence de connexion}
\begin{lstlisting}[language=mermaid]
sequenceDiagram
    actor U as Utilisateur
    participant F as Frontend React
    participant A as API Auth Django
    participant DB as SQLite
    U->>F: Saisir username/password
    F->>A: POST /api/auth/login/
    A->>DB: Verifier utilisateur
    DB-->>A: Utilisateur valide
    A-->>F: user + token
    F->>F: Stocker token localStorage
    F-->>U: Redirection dashboard
\end{lstlisting}
\end{figure}

\subsection{Diagramme de séquence : diagnostic image}
\begin{figure}[H]
\centering
\caption{Séquence d'analyse de maladie}
\begin{lstlisting}[language=mermaid]
sequenceDiagram
    actor U as Utilisateur
    participant F as Page Analysis
    participant API as API Analysis
    participant CNN as plant_disease_cnn.h5
    participant DB as SQLite
    U->>F: Selectionner image
    F->>API: POST /api/analysis/upload/
    API->>API: Valider type et taille
    API->>DB: Creer AnalysisLog
    API->>API: Resize 224x224, RGB, normalisation
    API->>CNN: Prediction
    CNN-->>API: Probabilites 38 classes
    API->>DB: Sauvegarder resultat JSON
    API-->>F: Top prediction + recommandations
    F-->>U: Afficher diagnostic
\end{lstlisting}
\end{figure}

\subsection{Diagramme d'activité : recommandation agricole}
\begin{figure}[H]
\centering
\caption{Activité du module de recommandation}
\begin{lstlisting}[language=mermaid]
flowchart TD
    A[Debut] --> B[Saisie localisation]
    B --> C[Geocodage Open-Meteo]
    C --> D{Meteo disponible ?}
    D -- Non --> E[Retour erreur 503]
    D -- Oui --> F[Charger donnees sol]
    F --> G{Mode manuel ?}
    G -- Oui --> H[Surcharger NPK, pH, temperature, humidite]
    G -- Non --> I[Utiliser donnees automatiques]
    H --> J[Construire DataFrame]
    I --> J
    J --> K[Predire avec RandomForest]
    K --> L[Calculer confiance et alternatives]
    L --> M[Generer conseils]
    M --> N[Retour JSON]
\end{lstlisting}
\end{figure}

\subsection{Diagramme de composants}
\begin{figure}[H]
\centering
\caption{Diagramme de composants}
\begin{lstlisting}[language=mermaid]
flowchart LR
    subgraph Frontend
        Pages[Pages React]
        Services[Services Axios]
        Context[AuthContext]
    end

    subgraph Backend
        Auth[authentication]
        Weather[weather]
        News[news]
        Analysis[analysis]
        Crops[crops]
        Chatbot[chatbot]
        Messaging[messaging]
    end

    subgraph Data
        SQLite[(db.sqlite3)]
        Media[(media files)]
        Models[(PKL/H5 models)]
    end

    Pages --> Services
    Context --> Services
    Services --> Auth
    Services --> Weather
    Services --> News
    Services --> Analysis
    Services --> Chatbot
    Services --> Messaging
    Pages --> Crops
    Auth --> SQLite
    Messaging --> SQLite
    Analysis --> SQLite
    Chatbot --> SQLite
    Analysis --> Media
    Analysis --> Models
    Crops --> Models
\end{lstlisting}
\end{figure}

\section{Architecture logicielle}
L'application utilise une architecture client-serveur de type trois tiers. Le frontend React constitue la couche présentation. Le backend Django REST constitue la couche métier et expose les APIs. La base SQLite, les fichiers médias et les fichiers modèles constituent la couche données.

\subsection{Architecture frontend}
Le frontend est organisé autour des dossiers :
\begin{itemize}[noitemsep]
  \item \texttt{pages} : pages routées par React Router ;
  \item \texttt{components} : composants réutilisables dont \texttt{Navbar}, \texttt{Footer} et \texttt{RecommendationForm} ;
  \item \texttt{services} : clients API Axios ;
  \item \texttt{contexts} : contexte d'authentification ;
  \item \texttt{styles} : feuilles CSS.
\end{itemize}

\subsection{Architecture backend}
Le backend est organisé en applications Django indépendantes :
\texttt{authentication}, \texttt{messaging}, \texttt{weather}, \texttt{news}, \texttt{analysis}, \texttt{chatbot}, \texttt{crops}. Cette séparation améliore la maintenabilité et limite le couplage.

\section{Base de données}
\subsection{MCD}
\begin{figure}[H]
\centering
\caption{MCD simplifié}
\begin{lstlisting}[language=mermaid]
erDiagram
    USER ||--o{ MESSAGE : sends
    USER ||--o{ MESSAGE : receives
    USER ||--o{ ANALYSIS_LOG : owns
    USER ||--o{ CHAT_MESSAGE : owns

    USER {
      int id
      string username
      string email
      string role
      string phone
      string avatar
    }
    MESSAGE {
      int id
      int sender_id
      int recipient_id
      string subject
      text content
      boolean is_read
    }
    ANALYSIS_LOG {
      int id
      int user_id
      string image
      json result
      float confidence
      float processing_time
    }
    CHAT_MESSAGE {
      int id
      int user_id
      text message
      text response
    }
\end{lstlisting}
\end{figure}

\subsection{Explication des entités}
\begin{itemize}[noitemsep]
  \item \texttt{User} : utilisateur personnalisé avec rôle, téléphone et avatar.
  \item \texttt{Message} : message interne entre deux utilisateurs.
  \item \texttt{AnalysisLog} : historique d'analyse image, résultat JSON et métriques d'inférence.
  \item \texttt{ChatMessage} : historique des échanges avec le chatbot.
\end{itemize}

\section{Gestion de projet SCRUM}
\subsection{Product backlog}
\begin{longtable}{p{1.2cm}p{8.5cm}p{2cm}p{2cm}}
\caption{Product backlog}\\
\toprule
\textbf{ID} & \textbf{User story} & \textbf{Priorité} & \textbf{Sprint} \\
\midrule
US1 & En tant que visiteur, je veux créer un compte. & Haute & S1 \\
US2 & En tant qu'utilisateur, je veux me connecter. & Haute & S1 \\
US3 & En tant qu'utilisateur, je veux consulter mon tableau de bord. & Haute & S1 \\
US4 & En tant qu'utilisateur, je veux consulter la météo agricole. & Haute & S2 \\
US5 & En tant qu'utilisateur, je veux consulter les actualités. & Moyenne & S2 \\
US6 & En tant qu'agriculteur, je veux obtenir une recommandation de culture. & Très haute & S2 \\
US7 & En tant qu'agriculteur, je veux personnaliser les valeurs du sol. & Haute & S3 \\
US8 & En tant qu'utilisateur, je veux diagnostiquer une maladie à partir d'une image. & Très haute & S3 \\
US9 & En tant qu'utilisateur, je veux recevoir des conseils de traitement. & Très haute & S3 \\
US10 & En tant qu'utilisateur, je veux discuter avec un chatbot. & Moyenne & S4 \\
US11 & En tant qu'utilisateur, je veux envoyer et recevoir des messages. & Moyenne & S4 \\
US12 & En tant qu'administrateur, je veux consulter un tableau de bord. & Moyenne & S4 \\
\bottomrule
\end{longtable}

\subsection{Planification des sprints}
\begin{table}[H]
\centering
\caption{Planification des sprints}
\begin{tabular}{p{2cm}p{5cm}p{7cm}}
\toprule
\textbf{Sprint} & \textbf{Objectif} & \textbf{Livrables} \\
\midrule
S1 & Socle applicatif & Structure React/Django, routage, authentification, dashboards. \\
S2 & Services contextuels & Météo, actualités, recommandation initiale de cultures. \\
S3 & Intelligence artificielle & Random Forest Tunisie, CNN maladies, upload image, recommandations. \\
S4 & Communication et consolidation & Messagerie, chatbot, dashboard admin, documentation et tests. \\
\bottomrule
\end{tabular}
\end{table}
